\section{StyleGAN and StyleGAN2}

Early GANs suffered from instability and mode collapse, which \citet{karras2018progressivegrowinggansimproved} addressed by 
Progressive GANs (PGGAN): training begins at low resolution and gradually adds 
layers as resolution increases \citep{Melnik_2024}.
Building on this, StyleGAN \citet{karras2019stylebasedgeneratorarchitecturegenerative}
proposed a novel style-based generator architecture that achieved state-of-the-art image quality and disentangled latent factors.
StyleGAN is a variant of GAN that builds on its framework by radically redesigning the generator.

Instead of feeding the latent code $\mathbf{z}$ directly into the generator, 
StyleGAN first maps $\mathbf{z}\sim\mathcal{N}(0,I)$ to an intermediate latent $\mathbf{w}=f(\mathbf{z})$ using an 8-layer fully-connected network \citep{karras2019stylebasedgeneratorarchitecturegenerative}.
This $\mathbf{w}$ lives in a new “style” space $\mathcal{W}$ that is learned to be less entangled than the original latent space. 
Learned affine transforms then convert $\mathbf{w}$ into per-layer style parameters $(y_s,y_b)$ that modulate each convolution via Adaptive Instance Normalization (AdaIN) \citep{karras2019stylebasedgeneratorarchitecturegenerative,karras2020analyzingimprovingimagequality}.
Concretely, for each feature map $x_i$ in a convolutional layer, StyleGAN applies:\[
\text{AdaIN}(x_i, y) = y_{s,i} \left( \frac{x_i - \mu(x_i)}{\sigma(x_i)} \right) + y_{b,i}
\], where $y_s$ and $y_b$ (scalars for channel $i$) are derived from $\mathbf{w}$.
 In practice, StyleGAN’s mapping network $f$ (an eight-layer MLP) produces $\mathbf{w}$, and then each convolutional layer is modulated by the corresponding $(y_s,y_b)$ pair.
 This allows scale-specific control of features: coarse spatial styles 
 (e.g. pose or overall shape) are set by early layers, while fine-grained details (e.g. color, texture) are set by later layers \cite{Melnik_2024,karras2019stylebasedgeneratorarchitecturegenerative}.
Notably, the mapping network and style modulation enable style-mixing: two latent codes can be injected at different depths, blending coarse features from one code with fine details from another \citep{Melnik_2024}.

Importantly, StyleGAN replaces the standard random input with a learned constant 4x4x512 tensor, and all variation comes from the style vectors and added noise. 
At each convolutional layer, StyleGAN also injects a random noise map (one per pixel) which is added to the features with learned per-channel weights.
This stochastic noise produces high-frequency variation (e.g. hair wisps, freckles) without affecting global attributes. 
Together, these design choices yield unsupervised disentanglement of high-level attributes from stochastic detail: \citet{karras2019stylebasedgeneratorarchitecturegenerative} observe that their style-based generator 
“leads to an automatically learned, unsupervised separation of high-level attributes (e.g., pose and identity) and stochastic variation in the generated images”.
In experiments, StyleGAN achieves state-of-the-art image fidelity (low Fréchet Inception Distance) and more linear, less entangled latent factors than conventional GANs.


The architectural differences from a standard GAN can be summarized as:

\begin{itemize}
  \item Intermediate latent space and mapping network: StyleGAN learns a mapping $f: \mathbf{z}\to \mathbf{w}$ (eight-layer MLP) so that $\mathbf{w}$ in space $\mathcal{W}$ is disentangled
  \item Style modulation via AdaIN: Style vectors from $\mathbf{w}$ provide per-layer \\ scale/bias that modulate each convolution by AdaIN, as in the formula above.
  \item Noise injection: Learned Gaussian noise is added to each layer to control stochastic details
  \item Learned constant input: The generator begins from a fixed learned 4x4 tensor instead of random noise
  \item Progressive growing: StyleGAN adopts the progressive training strategy from ProGAN [Karras et al., 2018], starting at 4x4 resolution and gradually adding layers to reach high resolution \citep{karras2018progressivegrowinggansimproved}.
\end{itemize}



StyleGAN2 \citet{karras2020analyzingimprovingimagequality} builds directly on StyleGAN but redesigns the generator to eliminate artifacts and boost quality.
The key insight is that the original AdaIN normalization can destroy meaningful signal: normalizing each feature map breaks global correlations needed for coherent details.
Therefore, StyleGAN2 “bakes” the style modulation into the convolution weights and replaces the per-feature normalization with a weight demodulation operation.
In practice, this works as follows: let $w_{ijk}$ be the convolution weight connecting input channel $i$ to output channel $j$ (at spatial position $k$).
Given a style scale $s_i$ (from $\mathbf{w}$) for channel $i$, StyleGAN2 first multiplies $w_{ijk}$ by $s_i$ (modulation), yielding $w'_{ijk}=s_i w_{ijk}$.
It then computes the output standard deviation $\sigma_j$ of the convolution and demodulates by dividing by $\sigma_j$:
\[
w''_{i,s,j,k} = \frac{w'_{ijk}}{\sqrt{\sum_{i'} (w'_{ijk})^2 + \epsilon}}
\]. This eliminates the need for explicit instance normalization after the convolution. 
The result is that each convolutional layer can apply the style $s_i$ directly to its weights, 
with the demodulation automatically preserving unit variance. This modification removes the “blob” artifacts common in StyleGAN outputs and leads to crisper, more natural images.

StyleGAN2 also revamps the training scheme. The progressive growing paradigm is replaced with a fixed architecture: the generator is trained 
at full resolution from the start using skip connections and residual blocks that connect multi-scale features (as in MSG-GAN or U-Net architectures).
By preserving multi-scale signals through the network rather than gradually adding layers, StyleGAN2 achieves stable training without the overhead of layer-by-layer growth.
Importantly, StyleGAN2 retains the style mapping network and noise injections, but applies them within this revised generator. 
Additionally, \citet{karras2020analyzingimprovingimagequality} introduce a path length regularization on the mapping network, which encourages equal-distance moves in $\mathcal{W}$ to produce equal perceptual changes in the image.
This regularizer makes the latent-to-image mapping smoother and more invertible, further improving disentanglement and editability.

The architectural improvements in StyleGAN2 can be summarized as:

\begin{itemize}
    \item Weight demodulation instead of AdaIN: Styles modulate convolution \\weights and are normalized in-weight, removing all explicit feature-map normalization.
    \item No instance normalization layers: By design, the generator no longer uses AdaIN or other instance/batch norms; all normalization is implicit in weight demodulation
    \item Elimination of progressive growing: The network is trained at fixed full resolution, using skip/residual connections to preserve stability
    \item Path length regularization: A regularizer on the mapping network encourages a smooth, linearized latent space.
\end{itemize}

These changes yield a substantial quality gain. 
For example, \citet{karras2019stylebasedgeneratorarchitecturegenerative} report that fully updated StyleGAN2 
(no progressive growing, weight demodulation, path regularization) achieves a Fréchet Inception Distance of 
about 3.3 on FFHQ, compared to ~4.4 for the original StyleGAN baseline.

While StyleGAN and StyleGAN2 achieve unprecedented image realism, they have known limitations. 
First, the generated images and editing capabilities are confined to the training distribution. 
StyleGAN's latent control only applies to images the network can produce; arbitrary real-world 
photos must first be inverted into $W$ (GAN inversion) before editing \citep{bermano2022stateoftheartarchitecturemethodsapplications}.
Even then, inversion may not perfectly capture identity. 
Second, the models excel on structured domains like faces or bedrooms but struggle with unstructured or complex scenes. 
As \citet{bermano2022stateoftheartarchitecturemethodsapplications} observed, StyleGAN “struggles with domains that do not exhibit strong structure”.
Third, bias is a concern: like all data-driven GANs, 
StyleGAN can replicate demographic biases in the training set, raising fairness issues. 
Finally, the architecture is computationally heavy and complex; 
training StyleGAN2 on high-resolution data demands extensive GPU resources and careful tuning.
